\chapter{Centralized Multi-Agent Theorem Proving}\label{chap:centralized}

\newcommand{\proves}{\vdash}




With the motivation of modelling real-world mathematical development in mind, we first introduce a formalism of the core structure of mathematical tasks that will be used throughout this thesis. Inspired by proof-theoretic and type-theoretic semantics, we represent mathematical statements as formulas within a formal language, and the process of proving these statements as the resolution of proof obligations (sequents) that are defined by these formulas. We then introduce a library structure that models the evolving state of mathematical knowledge across multiple agents, each with private information.

Then with this formalism in place, we introduce the centralized multi-agent theorem proving problem, a POMDP in which a single centralized planner allocates compute across proof/conjecture/publication actions for multiple agents with private information. We study this problem under the simplified environment of perfect public information, and study its complexity and approximation-theoretic properties by relating it to stochastic online scheduling problems with precedence constraints. We show that the problem is intractable in general, but identify special cases that are tractable and/or admit good approximations. We also discuss the limitations of this centralized formulation in the full environment of private information, and motivate the need for a decentralized formulation in which coordination is mediated by a market mechanism, which we study in the next chapter.


\subsection{Background}

\subsubsection{Syntax}

We consider the \textit{formula} $\phi$ as the core syntactic primitive of (formal) mathematics. Namely, we represent mathematical theorem statements syntactically as a formula (i.e. object of type \texttt{Prop}) within some formal language $\Sigma$, such as Lean, Rocq, etc. We define $\mathcal{F}$ as the collection of all such formulas, and equip it with an involution $\lnot : \mathcal{F}\rightarrow \mathcal{F}$; namely, $\lnot(\lnot\varphi)=\varphi$ for all $\varphi\in\mathcal{F}$. Note that we will refer to formulas as ``nodes'' interchangably due to a natural graphical structure that will be discussed later.

\subsubsection{Semantics}

We now aim to represent the semantic meaning of a theorem statement within a particular context concisely. Towards this, we consider the \textit{sequent}, given as an element of $\mathcal{S} = \mathcal{P}_{fin}(\mathcal{F}) \times \mathcal{F} $ and stylized as $[\Gamma \proves \phi]$ (read: ``Gamma \textit{proves} phi''). Namely, such a sequent represents a \textit{proof obligation}, i.e. the notion of proving a theorem statement $\phi$ (the \textit{consequent}) within the \textit{antecedent} context $\Gamma$. Such a proof obligation may be \textit{resolved}, meaning proven.

\subsubsection{Libraries}

We assume the existence of a universal clock $t \in \mathbb N$ and set of agents $\mathcal N$, and use it to define a library $\mathcal{L}_\bullet = \{\mathcal{L}_t\}_{t \in \mathbb{N}}$ as a collection of discrete library states that evolves over time. At every discrete timestep $t \in \mathbb{N}$, each library state is given as a tuple 
\[\mathcal{L}_t = (\mathcal{C}_t,\mathcal{RS}_t, \delta_t, \Theta_t) \in \mathscr{L}\]
where we partition $\mathcal{F} = \mathcal{C}_t \cup \mathcal{G}_t, \mathcal{C}_t \cap \mathcal{G}_t = \varnothing$ into \textit{concrete} and \textit{ghost} nodes. Intuitively, concrete nodes represent theorem statements that have been stated or conjectured, and ghost nodes represent all hypothetical, possible theorem statements that are yet to be conjectured or stated. $\mathcal{C}_t$ naturally induces a semantic analogue: let $\mathcal{CS}_t = \mathcal{P}_{fin}(\mathcal{C}_t) \times \mathcal{C}_t$ be the set of all \textit{concrete sequents}, which contains a subset $\mathcal{RS}_t \subseteq \mathcal{CS}_t$ of \textit{resolved sequents}, which represent the concrete nodes that have been proven (within a particular antecedent context, of course). 


Mathematically, a proof of a particular theorem in some context often explicitly depends on other lemmas within this context. In an effort to model this phenomenon, define $\delta_t : \mathcal{RS}_t \to \mathcal{P}_{fin}(\mathcal{F})$ as the dependency mapping.
Note that in mathematics, there may exist many proofs of the same theorem with different dependencies for each, and one may find new proofs of the same theorem that are also valid. Thus, in the most general formulation, we only require that resolved sequents must remain resolved ($\mathcal{RS}_{t} \subseteq \mathcal{RS}_{t+1}$), though the underlying ``proof'' (moreso, the set of dependencies of this ``proof'') can change.

And lastly, we define $\Theta_t : \mathcal{RS}_t\to\mathbb N^+ \times \mathcal N$ to record the (discrete) solve-time associated to each resolved sequent, and the agent responsible for the resolution. 


So, in other words, the library state $\mathcal{L}_t$ represents the collection of all stated theorem statements, the collection of proof witnesses within particular contexts for a subset of these statements, and the dependencies and ancillary metadata of these witnesses.

\subsubsection{Library Invariants and Properties}

Naturally, we impose a monotonicity invariant for libraries: one may not ``unconjecture'' or ``unprove'' a theorem. Formally, given a library $\mathcal{L}_\bullet = \{\mathcal{L}_t\}_{t \in \mathbb{N}}=\{(\mathcal{C}_t,\mathcal{RS}_t,\delta_t,\Theta_t)\}_{t \in \mathbb{N}}$, we require that for all $t \in \mathbb{N}$:
\[\mathcal{C}_{t} \subseteq \mathcal{C}_{t+1} \quad \text{and} \quad \mathcal{RS}_{t} \subseteq \mathcal{RS}_{t+1}\]

Additionally, we require that if one conjectures a theorem, they must also have automatically conjectured its converse; in other words:
\[\phi \in \mathcal{C}_t \iff \neg\phi \in \mathcal{C}_t\]

Similarly, one may not prove both a theorem and its converse true: 
\[[\Gamma \proves \phi] \in \mathcal{RS}_t \implies [\Gamma \proves \neg \phi] \notin \mathcal{RS}_t\]

Semantically, we note reflexivity holds: $\phi \in \Gamma \implies [\Gamma \proves \phi] \in \mathcal{RS}_t$. And for the sake of consistency, we require that for any $[\Gamma \proves \phi] \in \mathcal{CS}_t$, $\psi \in \Gamma \implies \neg\psi \notin \Gamma$ (although we do not impose any full propositional semantics or richer logical theory behind these abstract formulas, preferring to leave them fully symbolic objects with various resolution and concreteness states, allowing inconsistent antecedents in our witnesses limits the scalability of our approach to deployment as a strategy module for multiagent theorem proving in a fully-defined formal language and type theory). 

Note that this last condition implicitly redefines
\[\mathcal{CS}_t = \mathcal{P}_{fin}(\mathcal{C}_t)/\sim \times \mathcal{C}_t \] where $\phi\sim\neg\phi$ is the equivalence relation generated by the negation operation. So, due to this consistency condition, we note that $\mathcal{RS}_t$ is upwards-closed, in that if $[\Gamma \proves \phi] \in \mathcal{RS}_t$, then for all $\Gamma' \in \mathcal{P}_{fin}(\mathcal{C}_t)$, $\Gamma \subseteq \Gamma' \implies [\Gamma'\proves\phi]\in \mathcal{RS}_t$.

Additionally, considering the dependency mapping, we note that intuitively $\delta_t$ must satisfy the property that $\delta_t([\Gamma \proves \phi]) \subseteq \Gamma$. Moreover, one may use $\delta_t$ to induce a directed graph $D(\delta_t) = (\mathcal{CS}_t,\mathcal{E}_t)$, where 
\[\mathcal{E}_t = \Big\{([\Gamma\proves \phi], [\Gamma' \proves \psi]) : \phi \in \delta_t([\Gamma' \proves \psi]), \Gamma\subseteq \Gamma'\Big\}\]
We require that this induced directed graph $D(\delta_t)$ of $\delta_t$ must be acyclic.

\subsubsection{Axioms and Fully Resolved Sequents}

As an aside, we additionally note that although $[\Gamma \proves \phi] \in \mathcal{RS}_t$ may be marked resolved, that does not mean that $\phi$ has a ``error''-free proof within the current library state as its proof may depend on some still-unresolved lemmas. More concretely, we may define $\mathcal{FRS}_t \subseteq \mathcal{RS}_t$ as the collection of \textit{fully resolved sequents}, which is defined inductively such that \textit{axioms} $[\proves \phi] \in \mathcal{FRS}_t$ and $[\Gamma \proves \phi] \in \mathcal{FRS}_t$ if and only if for all $\phi' \in \delta_t([\Gamma \proves \phi])$, there exists $\Gamma'\subseteq \Gamma$ such that $[\Gamma' \proves \phi'] \in \mathcal{FRS}_t$.


\subsubsection{Private and Shared Libraries}
In the multiagent setting, each agent $\alpha\in\mathcal{N}$ maintains a \textit{private library} $\mathcal{L}_\bullet^\alpha = \{\mathcal{L}^\alpha_t\}_{t\in\mathbb{N}}$ with
\[
\mathcal{L}^\alpha_t = (\mathcal{C}^\alpha_t,\mathcal{RS}^\alpha_t,\delta^\alpha_t,\Theta_t^\alpha),
\]
and the environment maintains a \textit{shared (public) library} $\mathcal{L}_\bullet^0 = \{\mathcal{L}^0_t\}_{t\in\mathbb{N}}$ with
\[
\mathcal{L}^0_t = (\mathcal{C}^0_t,\mathcal{RS}^0_t,\delta^0_t, \Theta_t^0).
\]

We additionally enforce the following inclusion invariant for all $t \in \mathbb{N}, \alpha \in \mathcal{N}$:
\[
\mathcal{L}^0_t \subseteq \mathcal{L}^\alpha_t
\]
meaning $\mathcal{C}^0_t\subseteq \mathcal{C}^\alpha_t$, $\mathcal{RS}^0_t\subseteq \mathcal{RS}^\alpha_t$, and for all $s\in \mathcal{RS}_t^0, \delta^0_t(s) = \delta^\alpha_t(s)$ and $\Theta^0_t(s)=\Theta^\alpha_t(s)$.


\subsection{The Centralized MATP Problem}

We now introduce the Centralized Multi-Agent Theorem Proving (MATP) problem, a POMDP in which a single centralized planner allocates compute across proof/conjecture/publication actions for multiple agents with private information.

\subsubsection{Background: POMDP}
We first recall the standard Markov Decision Process (MDP) formalism, which models a single agent sequential decision making problem with full state observability.

\begin{definition}[Markov Decision Process]
A discrete-time Markov decision process is a tuple
\[
\mathcal{M} = (\mathsf{S},\;\mathsf{A},\;T,\;T^0,\;c),
\]
where $\mathsf{S}$ is a measurable state space, $\mathsf{A}$ is a measurable action space, $T(\cdot\mid s,a)$ is a transition kernel over $\mathsf{S}$, $T^0$ is an initial state distribution over $\mathsf{S}$, and $c:\mathsf{S}\times\mathsf{A}\to\mathbb{R}$ is a per-step cost function. At time $t=0$, nature draws $s_0\sim T^0$. At each subsequent time $t\ge 0$, the controller selects $a_t$, incurs cost $c(s_t,a_t)$, and the environment transitions to $s_{t+1}\sim T(\cdot\mid s_t,a_t)$.
\end{definition}

One may generalize the MDP formalism to include partial observability, which is more appropriate for modelling the multiagent theorem proving problem in its full generality, as each agent has private information in the form of their private library state. The natural formulation of this problem is as a Partially Observable Markov Decision Process (POMDP), in which the agent does not directly observe the underlying state, but rather receives observations that are generated from the underlying state according to some observation function.

\begin{definition}[POMDP]
A discrete-time partially observable Markov decision process is a tuple
\[
\mathcal{P} = (\mathsf{S},\; \mathsf{A},\; \mathsf{O},\; T,\; T^0,\; Z,\; c),
\]
for a state space $\mathsf{S}$, a measurable space whose elements encode the full (hidden) state of the environment; an action space $\mathsf{A}$, a measurable space of actions available to the controller; an observation space $\mathsf{O}$, a measurable space of observations received by the controller; a transition kernel $T(\cdot \mid s, a)$ over $\mathsf{S}$ specifying the distribution of the next state given a current state and action; an initial state distribution $T^0$ over $\mathsf{S}$; an observation kernel $Z(\cdot \mid s)$ over $\mathsf{O}$ specifying the distribution of the observation received by the controller when the environment is in state $s$; and a per-step cost function $c : \mathsf{S} \times \mathsf{A} \to \mathbb{R}$.
\end{definition}
 
The process evolves as follows. At time $t = 0$, nature draws a latent state
$s_0 \sim T^0$ and the controller receives an observation $o_0 \sim Z(\cdot \mid s_0)$.
At each subsequent time $t \ge 0$, the controller selects an action $a_t$, the environment
transitions to a new state $s_{t+1} \sim T(\cdot \mid s_t, a_t)$, and the controller receives
a new observation $o_{t+1} \sim Z(\cdot \mid s_{t+1})$.
 
The controller's observation history at time $t$ is the tuple
\[
h_t = (o_0,\, a_0,\, o_1,\, a_1,\, \ldots,\, a_{t-1},\, o_t),
\]
which constitutes the controller's complete information record. 
 
\begin{definition}[Policy]
A randomized nonanticipative policy is a family of stochastic kernels
$\pi = (\pi_t)_{t \ge 0}$ with $\pi_t : \mathsf{H}_t \to \mathcal{P}(\mathsf{A})$, where $\mathsf{H}_t = (\mathsf{O}\times\mathsf{A})^t\times\mathsf{O}$ is the space of observation histories of length $t$, such that $a_t \sim \pi_t(\cdot \mid h_t)$ depends only on the observation history up to time $t$. Equivalently, the action process $(a_t)_{t \ge 0}$ is adapted to the observation filtration $\mathcal{F}_t^{\mathrm{obs}} = \sigma(o_0, a_0, \ldots, a_{t-1}, o_t)$.
\end{definition}
 
\begin{remark}[Belief state]
By standard POMDP theory, the posterior belief
$b_t = \mathrm{Law}(s_t \mid h_t) \in \Delta(\mathsf{S})$ is a sufficient statistic of
the observation history for optimal control: there exists an optimal policy whose action
at time $t$ depends on $h_t$ only through $b_t$. One may equivalently reformulate the
POMDP as a fully observed MDP on the belief space $\mathcal{P}(\mathsf{S})$, the \emph{belief MDP}, in which the state is $b_t$ and the transition is given by Bayesian updating.
\end{remark}

To solve a POMDP means to find a policy $\pi$ that minimizes the expected total (discounted) cost:
\[
J(\pi) = \mathbb{E}^{\pi}\!\left[\sum_{t=0}^{\infty}\gamma^t\, c(s_t,a_t)\;\middle|\; s_0\sim T^0\right]
\]
for a discount factor $\gamma\in[0,1)$, or the expected total cost under a finite horizon $H$:
\[
J_H(\pi)=\mathbb{E}^\pi\!\left[\sum_{t=0}^{H-1}c(s_t,a_t)\;\middle|\;s_0\sim T^0\right].
\]

\subsubsection{Background: Action Primitives}
We now specify the action primitives that the centralized planner may assign to individual agents. These primitives are standalone kernels and update maps that will be assembled into the POMDP transition structure.
 
\paragraph{Discrete kernels.}
All stochastic primitives in this section have finite or countable output spaces, and may be treated as
probability mass functions (pmfs). Concretely, for an input space $X$ and discrete output space $Y$,
a stochastic kernel is a map $K:Y\times X\to[0,1]$ such that $\sum_{y\in Y}K(y\mid x)=1$ for all $x\in X$.
 
\paragraph{Proof Kernel.}
 
Fix an agent $\alpha\in\mathcal{N}$. The proof kernel models allocating a time budget to attempt to resolve a target sequent. Let the input space be
\[
X_{\mathrm{proof}} = \mathscr{L}\times\mathcal{S}\times \mathbb{N}^+,
\]
where the components are a library state, a target sequent $s=[\Gamma\proves\phi]$, and a discrete time budget $\tau$. Let the output space be
\[
Y_{\mathrm{proof}} = \{(0,\bot)\}\;\cup\;\{(1,d): d\in\mathcal{P}_{\mathrm{fin}}(\mathcal{F})\},
\]
where $(0,\bot)$ denotes failure, and $(1,d)$ denotes success together with a dependency set $d$. The proof kernel for agent $\alpha$ is a pmf
\[\mathcal{K}^{\alpha}_{\mathrm{proof}} : Y_{\mathrm{proof}}\times X_{\mathrm{proof}} \to [0,1].\]
On admissible inputs $(\mathcal{L},[\Gamma\proves\phi],\tau)$, a successful output $(1,d)$ satisfies $d\subseteq\Gamma$.
 
\paragraph{Conjecture Kernel.}
 
Fix an agent $\alpha\in\mathcal{N}$. The conjecture kernel models allocating a time budget to propose a new formula, anchored at an existing sequent. Let the input space be
\[
X_{\mathrm{conj}} =\mathscr{L}\times \mathcal{S}\times \mathbb{N}^+,
\]
and the output space be
\[
Y_{\mathrm{conj}} = \{(0,\bot)\}\;\cup\;\{(1,\psi): \psi\in\mathcal{F}\}.
\]
The conjecture kernel for agent $\alpha$ is a pmf
$\mathcal{K}^{\alpha}_{\mathrm{conj}} : Y_{\mathrm{conj}}\times X_{\mathrm{conj}} \to [0,1]$.
Given admissible input $(\mathcal{L}, [\Gamma \proves \phi], \tau)$ with $\mathcal{L} = (\mathcal{C},\mathcal{RS},\delta,\Theta)$ and a successful output $(1,\psi)$, we enforce $\psi \in \mathcal{F} \setminus \mathcal{C}$.
 
\paragraph{Dependency Closure for Publication.}
 
Publication is a deterministic primitive that copies a closure-completed subset of an agent's private library
into the shared library. Fix agent $\alpha$ and time $t$, and suppose $\delta_t^\alpha$ induces the DAG
$D(\delta_t^\alpha)=(\mathcal{CS}_t^\alpha,\mathcal{E}_t^\alpha)$. For any set of target resolved sequents $R\subseteq \mathcal{RS}_t^\alpha$, define its \emph{dependency closure}:
\begin{align*}
\mathrm{cl}^{\alpha}_t(R) = 
\Big\{\, &s\in \mathcal{RS}_t^\alpha : \exists\, r\in R \text{ such that there}\\& \text{is a directed path } s\to r
\text{ in } D(\delta_t^\alpha)\,\Big\}
\end{align*}
 
A publication payload is a pair $U_t^\alpha=(C,R)$ with $C\subseteq \mathcal{C}_t^\alpha,\; R\subseteq \mathcal{RS}_t^\alpha$. Its closure-completion is $\overline{U_t^\alpha}=(\overline{C},\overline{R})$ where $\overline{R}=\mathrm{cl}^\alpha_t(R)$ and
\[\overline{C}= C \;\cup\; \bigcup_{[\Gamma\proves\phi]\in \overline{R}}(\Gamma\cup\{\phi\}).\]
 
\paragraph{Query Model.}
 
Each agent $\alpha$ maintains a stateful query model $Q_t^\alpha \in \mathcal{Q}$ that estimates prover performance. A query readout is a map
$q^\alpha : \mathcal{Q}\times \mathcal{S} \to [0,1]\times \mathbb{N}$, returning $(\hat p_t^\alpha(s),\hat \tau_t^\alpha(s)) = q^\alpha(Q_t^\alpha,s)$, where $\hat p_t^\alpha(s)$ estimates the probability that sequent $s$ is resolvable by agent $\alpha$, and $\hat\tau_t^\alpha(s)$ estimates the expected time-to-solution.
 
The query model is updated via an operator $\Delta^\alpha : \mathcal{Q}\times \mathcal{D}(\mathcal{Q}) \to \mathcal{Q}$, where $\mathcal{D}(\mathcal{Q}) = \mathcal{S} \times (\mathbb{N}^+ \times \mathcal{N})$ is the space of query-training data. We enforce commutativity and associativity of $\Delta^\alpha$ so that batch updates are well-defined and order-independent.
 
 
\subsubsection{Definition}
 
Fix a set of agents $\mathcal{N}$ and a target set of sequents $\Phi^*\subseteq \mathcal{S}$. The \emph{Centralized Multi-Agent Theorem Proving} (MATP) problem is specified by the POMDP
\[
\mathcal{P}_{\mathrm{MATP}} = (\mathsf{S},\;\mathsf{A},\;\mathsf{O},\;T,\;T^0,\;Z,\;c)
\]
defined as follows.
 
 
\paragraph{State Space $\mathsf{S}$.}
 
A state at time $t$ is a tuple
\begin{align*}
s_t =
\Big(
&\{\mathcal{L}^\alpha_t\}_{\alpha\in\mathcal{N}},\;
\mathcal{L}^0_t,\;
\{J^\alpha_t\}_{\alpha\in\mathcal{N}},\;
\{Q^\alpha_t\}_{\alpha\in\mathcal{N}},
\\&
\{\Omega^\alpha_t\}_{\alpha\in\mathcal{N}},\;
\{U_{t}^\alpha\}_{\alpha\in\mathcal{N}}
\Big)\in\mathsf{S},
\end{align*}
where:
\begin{itemize}
\item $\mathcal{L}^\alpha_t\in\mathscr{L}$ is agent $\alpha$'s private library and
$\mathcal{L}^0_t\in\mathscr{L}$ is the shared public library, subject to the inclusion invariant $\mathcal{L}^0_t \subseteq \mathcal{L}^\alpha_t$ for all $\alpha\in\mathcal{N}$.
\item $J^\alpha_t$ is the \emph{durative job descriptor}:
\[
J^\alpha_t \in \{\bot\}\ \cup\ \big(\{\mathrm{Proof}, \mathrm{Conj}\}\times\mathcal{S}\times\mathbb{N}^+\times\mathbb{N}^+\big),
\]
with $\bot$ meaning idle and $(\cdot,s,\tau_{\mathrm{rem}},\tau_{\mathrm{eff}})$ encoding a job of type $\mathrm{Proof}$ or $\mathrm{Conj}$ on target sequent $s$ with $\tau_{\mathrm{rem}}$ steps remaining in the current run and $\tau_{\mathrm{eff}}$ the effective cumulative budget used to evaluate the kernel upon completion.
\item $Q^\alpha_t\in\mathcal{Q}$ is agent $\alpha$'s query model state.
\item $\Omega^\alpha_t \in \{\bot\}\ \sqcup\ (\mathcal{S}\times[0,1]\times\mathbb{N})$ is agent $\alpha$'s last query response, either empty $\bot$ or a triple $(s,\hat p,\hat \tau)$.
\item $U^\alpha_{t} : \mathcal{S} \times \{\mathrm{proof},\mathrm{conj}\} \to \mathbb{N}$ records the total time already spent by agent $\alpha$ on proof or conjecture attempts for each sequent, with $U^\alpha_0 \equiv 0$.
\end{itemize}
We take $\mathsf{S}$ to be the set of all such tuples satisfying the constituent objects' invariants.
 
 
\paragraph{Action Space $\mathsf{A}$ and Admissibility.}
 
The centralized planner selects a \emph{joint action} $a_t = (a^\alpha_t)_{\alpha\in\mathcal{N}} \in \mathsf{A}$, where $\mathsf{A}= \prod_{\alpha\in\mathcal{N}}\mathcal{A}^\alpha$ and each per-agent action space is
\begin{align*}
\mathcal{A}^\alpha
=\;&
\{\mathrm{Prove}\}\times(\mathcal{S}\times\mathbb{N}^+)
\\&\sqcup\;
\{\mathrm{Conj}\}\times(\mathcal{S}\times\mathbb{N}^+)
\\&\sqcup\;
\{\mathrm{Pub}\}\times\big(\mathcal{P}_{\mathrm{fin}}(\mathcal{F})\times \mathcal{P}_{\mathrm{fin}}(\mathcal{S})\big)
\\&\sqcup\;
\{\mathrm{Qry}\}\times \mathcal{S}
\\&\sqcup\;\{\bot\}.
\end{align*}
 
Define the admissibility predicate $\mathrm{Adm}:\mathsf{S}\times\mathsf{A}\to\{0,1\}$ componentwise: a joint action $a=(a^\alpha)_{\alpha\in\mathcal{N}}$ is admissible at state $s$ if and only if $\mathrm{Adm}^\alpha(s,a^\alpha)=1$ for all $\alpha\in\mathcal{N}$, where:
\begin{itemize}
    \item If $J^\alpha_t\neq\bot$, the agent is occupied and forced to no-op:
    \[
    \mathrm{Adm}^\alpha(s_t,a^\alpha)=\mathbf{1}[a^\alpha=\bot].
    \]
    \item If $J^\alpha_t=\bot$, the per-action admissibility conditions are:
    \begin{itemize}
    \item $\mathrm{Adm}^\alpha(s_t,\bot)=1$ always.
    \item For $a^\alpha = (\mathrm{Prove},(s,\tau))$ or $(\mathrm{Conj},(s,\tau))$:
    $\mathrm{Adm}^\alpha(s_t,a^\alpha)=1 \iff s \in \mathcal{CS}_t^\alpha.$
    \item For $a^\alpha = (\mathrm{Pub},(C,R))$: $\mathrm{Adm}^\alpha(s_t,a^\alpha)=1 \iff C\subseteq\mathcal{C}_t^\alpha$ and $R\subseteq\mathcal{RS}_t^\alpha.$
    \item For $a^\alpha = (\mathrm{Qry},s)$: $\mathrm{Adm}^\alpha(s_t,a^\alpha)=1 \iff s \in \mathcal{CS}_t^\alpha.$
    \end{itemize}
\end{itemize}
 
\begin{remark}[Centralization and admissibility]
\label{rem:centralized-admissibility}
Note the admissibility predicate for agent $\alpha$ depends on $\alpha$'s private library $\mathcal{L}^\alpha_t$, which the centralized planner may not fully observe. Thus, the planner must reason about which actions are admissible for each agent based on its belief over the hidden state. 
\end{remark}
 
 
\paragraph{Observation Space $\mathsf{O}$ and Observation Kernel $Z$.}
 
The centralized planner observes the public projection of the state:
\[
o_t =
\big(\mathcal{L}^0_t,\;\{\Omega^\alpha_t\}_{\alpha\in\mathcal{N}}\big)
\in\mathsf{O}.
\]
That is, the planner observes the shared library and the most recent query response from each agent. In particular, the private libraries $\{\mathcal{L}^\alpha_t\setminus\mathcal{L}^0_t\}_{\alpha\in\mathcal{N}}$, internal job states, and query model parameters are \emph{hidden}.
 
The observation kernel is deterministic:
\[
Z(o\mid s)=\mathbf{1}\!\big[\,o=\mathrm{Obs}(s)\,\big],
\]
where $\mathrm{Obs}:\mathsf{S}\to\mathsf{O}$ computes the public projection as above.
 
\begin{remark}[Perfect public information]
\label{rem:perfect-public-info}
In the special case where all libraries are public and publication is forced on resolution (i.e., $\mathcal{L}^\alpha_t = \mathcal{L}^0_t$ for all $\alpha\in\mathcal{N}$ and all $t$), the observation $o_t$ determines $s_t$ and the POMDP reduces to a fully observed MDP. We refer to this as the \emph{perfect public information} (PPI) setting and study it in Section~\ref{sec:complexity}.
\end{remark}
 
 
\paragraph{Transition Kernel $T$.}
 
We define $T(\cdot\mid s_t,a_t)$ by a three-step generative construction: a deterministic pre-update, a stochastic job-completion draw, and a deterministic post-update.
 
\subparagraph{Step 1: Deterministic Pre-Update.}
 
Given a state $s_t\in\mathsf{S}$ and admissible joint action $a_t\in\mathsf{A}$, define
$\tilde s_t = \mathrm{Pre}(s_t,a_t)$ by applying the following edits in order:
\begin{enumerate}
    \item \textbf{Decrement job timers.} For each $\alpha \in \mathcal{N}$: if $J^\alpha_t=\bot$, leave it unchanged; if $J^\alpha_t=(\cdot,s,\tau_{\mathrm{rem}},\tau_{\mathrm{eff}})$ with $\tau_{\mathrm{rem}}>0$, set
$\tau_{\mathrm{rem}}\leftarrow \tau_{\mathrm{rem}}-1$.
 
\item \textbf{Start new jobs.} For each idle agent $\alpha$ (i.e., $J^\alpha_t=\bot$):
\begin{itemize}
    \item If $a^\alpha_t=(\mathrm{Prove},(s,\tau))$, set
$J^\alpha \leftarrow (\mathrm{Proof},s,\tau,\tau+U^\alpha_t(s,\mathrm{proof}))$.
	\item If $a^\alpha_t=(\mathrm{Conj},(s,\tau))$, set
$J^\alpha \leftarrow (\mathrm{Conj},s,\tau,\tau+U^\alpha_t(s,\mathrm{conj}))$.
\end{itemize}
 
\item \textbf{Apply publications.} For each $\alpha\in\mathcal{N}$ with $a^\alpha_t=(\mathrm{Pub},(C^\alpha,R^\alpha))$: compute the closure-completion $\overline{(C^\alpha,R^\alpha)}=(\overline{C}^\alpha,\overline{R}^\alpha)$ and update:
\[\mathcal C^0 \leftarrow \mathcal C^0 \cup \textstyle\bigcup_{\alpha}\overline C^\alpha,
\qquad
\mathcal{RS}^0 \leftarrow \mathcal{RS}^0 \cup \textstyle\bigcup_{\alpha}\overline R^\alpha.\]
For each newly resolved $s \in \mathcal{RS}^0_{\mathrm{new}} = \mathcal{RS}^0 \setminus \mathcal{RS}^0_{\mathrm{old}}$, set $\delta^0(s) \leftarrow \delta^\alpha(s)$ and $\Theta^0(s) \leftarrow \Theta^\alpha(s)$.
Propagate to each agent $\beta$:
\[\mathcal C^\beta \leftarrow \mathcal C^\beta \cup (\mathcal C^0\setminus \mathcal C^0_{\mathrm{old}}),\quad
\mathcal{RS}^\beta \leftarrow \mathcal{RS}^\beta \cup \mathcal{RS}^0_{\mathrm{new}},\]
and for each $s\in \mathcal{RS}^0_{\mathrm{new}}$:
$\delta^\beta(s)\leftarrow\delta^0(s)$, $\Theta^\beta(s)\leftarrow\Theta^0(s)$, and $Q^\beta \leftarrow \Delta^\beta(Q^\beta,(s,\Theta^0(s)))$.
 
\item \textbf{Update query buffers.} For each $\alpha\in\mathcal{N}$: if $a^\alpha_t=(\mathrm{Qry},s)$, set
$\Omega^\alpha \leftarrow (s,\hat p,\max\{0,\hat\tau-U^\alpha_t(s,\mathrm{proof})\})$
where $(\hat p,\hat\tau)=q^\alpha(Q^\alpha,s)$; otherwise set $\Omega^\alpha\leftarrow \bot$.
\end{enumerate}
 
\subparagraph{Step 2: Stochastic Job Completion.}
 
Given $\tilde s_t$, define the sets of completing agents:
\[
    \mathsf{Comp}_{\mathrm{proof}}(\tilde s_t) = \{\alpha : J^\alpha(\tilde s_t)=(\mathrm{Proof},s,0,\tau)\text{ for some }s,\tau\},
\]
\[
\mathsf{Comp}_{\mathrm{conj}}(\tilde s_t) = \{\alpha : J^\alpha(\tilde s_t)=(\mathrm{Conj},s,0,\tau)\text{ for some }s,\tau\}.
\]
 
For each completing agent, we draw an independent outcome from the appropriate kernel. The joint outcome space is
\[Y(\tilde s_t)
=
\prod_{\alpha\in\mathsf{Comp}_{\mathrm{proof}}(\tilde s_t)} Y_{\mathrm{proof}}
\;\times\;
\prod_{\alpha\in\mathsf{Comp}_{\mathrm{conj}}(\tilde s_t)} Y_{\mathrm{conj}},\]
equipped with the product measure
\begin{align*}
\mathbb K_{\tilde s_t}(dy)
&=
\prod_{\alpha\in\mathsf{Comp}_{\mathrm{proof}}(\tilde s_t)}
\mathcal K^\alpha_{\mathrm{proof}}\!\big(dy^\alpha \mid \mathcal L^\alpha(\tilde s_t), s^\alpha_{\mathrm{proof}},\tau^\alpha_{\mathrm{proof}}\big)
\\&\quad\cdot
\prod_{\alpha\in\mathsf{Comp}_{\mathrm{conj}}(\tilde s_t)}
\mathcal K^\alpha_{\mathrm{conj}}\!\big(dy^\alpha \mid \mathcal L^\alpha(\tilde s_t), s^\alpha_{\mathrm{conj}},\tau^\alpha_{\mathrm{conj}}\big),
\end{align*}
where $(s^\alpha_{\cdot},\tau^\alpha_{\cdot})$ are read from $J^\alpha(\tilde s_t)$.
 
\subparagraph{Step 3: Deterministic Post-Update.}
 
Given $\tilde s_t$ and realized outcomes $y\in Y(\tilde s_t)$, define $s_{t+1} = \mathrm{Post}(\tilde s_t,y)$ by applying:
\begin{enumerate}
    \item \textbf{Resolve proof completions.} For each $\alpha\in\mathsf{Comp}_{\mathrm{proof}}(\tilde s_t)$ with outcome $y^\alpha = (b^\alpha, z^\alpha)$ and stored target/budget $(s^\alpha,\tau^\alpha)$: set $U^\alpha(s^\alpha,\mathrm{proof}) \leftarrow \tau^\alpha$. If $b^\alpha=1$ and $z^\alpha=d$:
\begin{itemize}
    \item $\mathcal{RS}^\alpha \leftarrow \mathcal{RS}^\alpha \cup \{s^\alpha\}$,
    \item $\delta^\alpha(s^\alpha) \leftarrow d$ \;(subject to $d\subseteq \Gamma$ and acyclicity),
    \item $\Theta^\alpha(s^\alpha)\leftarrow (\tau^\alpha,\alpha)$,
    \item $Q^\alpha \leftarrow \Delta^\alpha(Q^\alpha,(s^\alpha,\Theta^\alpha(s^\alpha)))$.
\end{itemize}
 
\item \textbf{Resolve conjecture completions.} For each $\alpha\in\mathsf{Comp}_{\mathrm{conj}}(\tilde s_t)$ with outcome $y^\alpha = (b^\alpha, z^\alpha)$: set $U^\alpha(s^\alpha,\mathrm{conj}) \leftarrow \tau^\alpha$. If $b^\alpha=1$ and $z^\alpha=\psi$:
$\mathcal C^\alpha \leftarrow \mathcal C^\alpha \cup \{\psi,\neg\psi\}$.
 
\item \textbf{Clear completed jobs.} For each $\alpha\in \mathsf{Comp}_{\mathrm{proof}}(\tilde s_t)\cup \mathsf{Comp}_{\mathrm{conj}}(\tilde s_t)$: set $J^\alpha \leftarrow \bot$.
\end{enumerate}
 
 
\paragraph{Full transition kernel.}
 
The transition kernel is the pushforward:
\[
T(U\mid s_t,a_t)
\;=\;
\int_{Y(\tilde s_t)}
\mathbf 1\!\big[\mathrm{Post}(\tilde s_t,y)\in U\big]\,
\mathbb K_{\tilde s_t}(dy),
\]
for measurable $U\subseteq \mathsf{S}$, where $\tilde s_t=\mathrm{Pre}(s_t,a_t)$.
 
 
\paragraph{Cost Function.}
 
Fix a target set $\Phi^*\subseteq \mathcal{CS}^0_0$ of sequents that the planner aims to fully resolve in the shared library. We define the per-step cost as
\[
c(s_t,a_t) = \sum_{s\in\Phi^*} \mathbf{1}\!\big[s\notin \mathcal{FRS}^0_t\big],
\]
the number of target sequents not yet fully resolved in the shared library at time $t$.
 
\begin{remark}[Cost variants]
Under this cost function, the finite-horizon objective $J_H(\pi) = \sum_{t=0}^{H-1}\mathbb{E}[c(s_t,a_t)]$ is minimized by policies that resolve target sequents as early as possible, yielding a natural connection to weighted completion time in scheduling theory. For a pure \emph{makespan} objective (minimizing the time at which all targets are resolved), one may instead consider the stopping-time formulation
\[
C_{\max}(\pi) = \inf\!\big\{t\ge 0 : \Phi^*\subseteq \mathcal{FRS}^0_t\big\}
\]
and seek to minimize $\mathbb{E}^\pi[C_{\max}(\pi)]$. The discounted formulation with $c$ as above interpolates between these via the discount factor $\gamma$.
\end{remark}


\paragraph{Initial distribution $T^0$.}

$T^0$ is a distribution over initial states $s_0$ specifying initial private and shared libraries (satisfying $\mathcal{L}^0_0\subseteq\mathcal{L}^\alpha_0$), idle job statuses ($J^\alpha_0=\bot$ for all $\alpha$), zero cumulative-effort maps ($U^\alpha_0\equiv 0$), initial query models $Q^\alpha_0\in\mathcal{Q}$, and empty query buffers ($\Omega^\alpha_0=\bot$), supported on states satisfying all invariants.


%%% Everything from \subsubsection{Example} onward %%%
%%% Replaces the existing content starting at that point %%%


\subsubsection{Example}\label{sec:matp-example}

We work through a simple instance to illustrate the centralized MATP formulation and the role of stochasticity and parallelism.

\paragraph{Setup.} Let $\mathcal{N}=\{\alpha_1,\alpha_2\}$ be two agents with identical proof kernels. The shared library $\mathcal{L}^0_0$ contains four concrete formulas $\{A,B_1,B_2,C\}$ (and their negations), with the single axiom $[\proves A]\in\mathcal{FRS}^0_0$. The planner's target is $\Phi^*=\{[A,B_1,B_2 \proves C]\}$. The (latent) dependency structure is:
\[
\delta^*([\Gamma \proves C]) = \{B_1,B_2\}, \qquad \delta^*([\Gamma \proves B_1]) = \varnothing, \qquad \delta^*([\Gamma \proves B_2]) = \varnothing.
\]
Thus $C$ depends on two independent lemmas $B_1$ and $B_2$, each provable from $A$. All libraries are public (PPI), so the POMDP reduces to an MDP. Each proof attempt takes $\tau=1$ step and succeeds independently with probability $p=0.5$.

\paragraph{Dependencies are latent.} At $t=0$, the planner knows the target $[A,B_1,B_2 \proves C]$ but does not know $\delta^*([\Gamma \proves C]) = \{B_1,B_2\}$. The planner may speculatively attempt $[A\proves B_1]$ or $[A\proves B_2]$, but it does not know these are \emph{needed} until $[\Gamma \proves C]$ is resolved and $\delta^*$ is revealed.


Suppose the planner first assigns both agents to attempt $[A,B_1,B_2 \proves C]$ (redundantly). If either succeeds at $t=1$ (as is expected), the planner learns $\delta^*([\Gamma \proves C]) = \{B_1,B_2\}$ and can assign each agent to one of $[A\proves B_1]$ and $[A\proves B_2]$ in parallel. The two independent lemmas are then resolved (in expectation) by $t=2$, and $C$ becomes fully resolved. Total expected makespan: $1/0.5 + 1/0.5 = 4$ (two serial phases, each with geometric waiting time $1/0.5$ when using redundant approaches with $m=2$ agents on a single task, giving effective success probability $1-(1-0.5)^2$).


Note that if the planner had perfect foresight about $\delta^*$, it could assign $\alpha_1\mapsto [A\proves B_1]$ and $\alpha_2\mapsto [A\proves B_2]$ at $t=0$, resolving both in parallel. It would then attempt $[A,B_1,B_2 \proves C]$ at $t=1$. This achieves makespan $\approx 1/0.5 + 1/0.5 = 2/0.5 = 4$ as well, so the latent dependency structure forces at least two serial phases regardless of foresight. This illustrates how the revelation DAG's depth (here, depth=$1$) imposes an irreducible serial bottleneck. Moreover, we see that the aforementioned policy is optimal compared to a perfect-foresight policy.


\subsection{Optimality of Immediate Publication}\label{sec:ppi-wlog}

The centralized MATP as defined allows private libraries: an agent's proof results initially reside in $\mathcal{L}^\alpha_t$ and enter the shared library $\mathcal{L}^0_t$ only upon an explicit $\mathrm{Pub}$ action. We now show that the PPI assumption ($\mathcal{L}^\alpha_t = \mathcal{L}^0_t$ for all $\alpha,t$) is without loss of generality for the centralized planner, in the sense that the optimal makespan under PPI is a lower bound on the optimal makespan in the full model.

\begin{theorem}[PPI as a relaxation]\label{thm:ppi-relaxation}
Let $\mathrm{OPT}$ denote the optimal expected makespan of the centralized MATP (with private libraries and explicit publication), and let $\mathrm{OPT}_{\mathrm{PPI}}$ denote the optimal expected makespan under the PPI assumption. Then:
\[
\mathrm{OPT}_{\mathrm{PPI}} \;\le\; \mathrm{OPT}.
\]
\end{theorem}

\begin{proof}
Under PPI, every proof completion is immediately visible in the shared library, namely, no publication action is required. Not only does that save a timestep and free computational time, but PPI also provides the planner strictly more information.

Under PPI, the planner observes the full state $s_t$ (since $o_t$ determines $s_t$ when all libraries coincide). In the non-PPI model, the planner observes only the public projection $o_t = (\mathcal{L}^0_t, \{\Omega^\alpha_t\}_\alpha)$ and must maintain beliefs over private library states. More information can only help (weakly) in a POMDP: the optimal value of the belief MDP is no worse than the optimal value of the partially observed problem.

Thus, every policy achievable in the non-PPI model can be simulated in the PPI model with equal or better makespan (the PPI model has more time and more information), so $\mathrm{OPT}_{\mathrm{PPI}} \le \mathrm{OPT}$.
\end{proof}

\begin{remark}[PPI as WLOG for lower bounds]
Theorem~\ref{thm:ppi-relaxation} means that any lower bound on $\mathrm{OPT}_{\mathrm{PPI}}$ is also a lower bound on $\mathrm{OPT}$, and any hardness result for the PPI problem implies hardness for the full problem. Conversely, any algorithm for the PPI problem yields an algorithm for the full problem by adding a publication subroutine (at the cost of additional timesteps for publication).
\end{remark}

\begin{remark}[When PPI is tight]\label{rem:ppi-tight}
The gap $\mathrm{OPT} - \mathrm{OPT}_{\mathrm{PPI}}$ is exactly the cost of publication in the full model. In instances where the number of publication events is small relative to the total computation (e.g., long proof attempts with infrequent publication), this gap is negligible. In the extreme case where all agents publish every resolved sequent as it completes and publication is instantaneous (zero-cost publication), the models coincide exactly.
\end{remark}

In light of Theorem~\ref{thm:ppi-relaxation}, we adopt the PPI assumption for the remainder of this chapter. Under PPI, the centralized MATP is a fully observed (stochastic) MDP: the planner's observation $o_t = (\mathcal{L}^0_t, \{\Omega^\alpha_t\}_\alpha)$ determines the full state $s_t$, publication is vacuous, and the planner selects joint actions with full knowledge of all agents' library states, job descriptors, and cumulative efforts.


\subsection{Complexity of the Centralized MATP}\label{sec:complexity}

We now establish the computational complexity of the centralized MATP under PPI, with and without conjecturing. Throughout, $m = |\mathcal{N}|$ is the number of agents.


\subsubsection{The Non-Conjecturing PPI MATP}\label{sec:nc-ppi}

We first isolate the scheduling-theoretic core by disabling conjecturing and imposing unique dependencies.

\begin{definition}[Non-Conjecturing PPI MATP]\label{def:nc-ppi-matp}
A \emph{Non-Conjecturing PPI MATP} instance is a centralized MATP instance under PPI, where the conjecture action is disabled, and $\mathcal{C}_0 = \mathcal{F}$: every formula is already concrete. The sequent space $\mathcal{CS}_0 = \mathcal{P}_{\mathrm{fin}}(\mathcal{F})/{\sim}\;\times\;\mathcal{F}$ is fixed for all time.
\end{definition}

Under these restrictions, the planner's decision at each timestep reduces to: for each idle agent, assign it to attempt a proof of some concrete sequent with some time budget, or leave it idle. The state evolves deterministically except for the stochastic outcomes of proof attempts.

\paragraph{Latent task set and revelation DAG.}

We recall that $\delta^*:\mathcal{S}\to\mathcal{P}_{\mathrm{fin}}(\mathcal{F})$ is a unique, latent deterministic function such that whenever any agent $\alpha$ successfully resolves $s$, the dependency set produced is $\delta^*(s)$. Namely, the planner does not know $\delta^*(s)$ until $s$ is resolved.

The unique-dependency assumption induces a deterministic (but latent) task set $\mathcal{T}^* = \mathrm{cl}(\Phi^*)$---the transitive dependency closure of the target set under $\delta^*$---and a \emph{revelation DAG} $G^{\mathrm{rev}} = (\mathcal{T}^*, E^{\mathrm{rev}})$ where $(s,s')\in E^{\mathrm{rev}}$ whenever $s'$ is induced by some $\psi\in\delta^*(s)$. A \emph{revelation chain} is a directed path in $G^{\mathrm{rev}}$; its \emph{depth} is the length of the longest such chain. The planner does not know $\mathcal{T}^*$ or $G^{\mathrm{rev}}$ at time $0$: they are revealed incrementally as proofs complete.

The key structural feature is that there is no \emph{operational} precedence: the planner may attempt any sequent at any time. The dependency structure constrains the \emph{completion criterion} (full resolution requires all transitive dependencies to be resolved), not the execution order. However, the latent revelation structure imposes an \emph{informational} bottleneck: the planner cannot learn that a sequent is needed until its parent in $G^{\mathrm{rev}}$ is resolved.


\subsubsection{Lower Bounds}\label{sec:lower-bounds}

\begin{definition}[Work and serial depth]\label{def:work-depth}
For a realization $\omega$ (determining all processing times), define:
\begin{itemize}
    \item The \emph{work}: $\displaystyle W(\omega) = \sum_{s\in\mathcal{T}^*} \min_{\alpha\in\mathcal{N}} T^\alpha_s(\omega)$.
    \item The \emph{serial depth}: $\displaystyle D(\omega) = \max_{\text{chain } s_0\to\cdots\to s_k \text{ in } G^{\mathrm{rev}}} \sum_{i=0}^{k} \min_{\alpha} T^\alpha_{s_i}(\omega)$.
\end{itemize}
\end{definition}

\begin{theorem}[Information-theoretic lower bounds]\label{thm:lower-bounds}
For any Non-Conjecturing PPI MATP instance and any admissible policy $\pi$:
\[
\mathbb{E}^\pi[C_{\max}] \;\ge\; \max\!\left\{\frac{\mathbb{E}[W]}{m},\;\; \mathbb{E}[D]\right\}.
\]
\end{theorem}

\begin{proof}
\emph{Work bound.} Each task $s\in\mathcal{T}^*$ must be resolved by some agent. With $m$ agents working in parallel, makespan $\ge W(\omega)/m$ for each $\omega$. Taking expectations gives the bound.

\emph{Serial depth bound.} Consider any revelation chain $s_0\to s_1\to\cdots\to s_k$ in $G^{\mathrm{rev}}$. Since $\delta^*$ is latent, the planner does not know $s_{i+1}\in\delta^*(s_i)$ until $s_i$ is resolved. Thus $s_{i+1}$ cannot be completed before time $C_{s_i} + \min_\alpha T^\alpha_{s_{i+1}}(\omega)$. Iterating, the chain requires time $\ge \sum_{i=0}^k \min_\alpha T^\alpha_{s_i}(\omega)$. Maximizing over chains and taking expectations gives $\mathbb{E}[C_{\max}] \ge \mathbb{E}[D]$.
\end{proof}

\begin{remark}[Speculation cannot beat the serial depth bound]
A speculative policy that resolves $s_{i+1}$ \emph{before} $s_i$ completes (guessing it will be needed) does not violate the bound. If $s_{i+1}$ has already been resolved speculatively when $s_i$ completes, then $\min_\alpha T^\alpha_{s_{i+1}} = 0$ in the chain, and the serial depth adjusts accordingly. Speculation can reduce the \emph{realized} serial depth but cannot reduce the \emph{ex ante} serial depth below $\mathbb{E}[D]$ without additional information about $\delta^*$.
\end{remark}


\subsubsection{Hardness}\label{sec:hardness}

\begin{theorem}[Complexity of the Non-Conjecturing PPI MATP]\label{thm:complexity}
\leavevmode
\begin{enumerate}
    \item \textbf{NP-hardness.} The decision problem---``given an instance $\mathcal{P}$ and threshold $T$, is $\mathrm{OPT}_{\mathrm{PPI}} \le T$?''---is NP-hard, even restricted to identical agents, deterministic proof kernels, and trivial dependencies ($\delta^*\equiv\varnothing$).
    
    \item \textbf{MDP membership.} Under PPI, the MATP is a fully observed MDP. The decision problem lies in PSPACE (by simulating the MDP in polynomial space).
    
    \item \textbf{State space exponentiality.} The MDP state space has cardinality $|\mathsf{S}| \ge 2^{|\mathcal{T}^*|}$, so exact solution via value iteration requires time exponential in the number of tasks.
\end{enumerate}
\end{theorem}

\begin{proof}
(1) We reduce from $P\|\, C_{\max}$ (makespan minimization on identical parallel machines), which is strongly NP-hard for variable $m$. Given an instance with $n$ jobs of processing times $p_1,\ldots,p_n$ on $m$ machines, construct a Non-Conjecturing PPI MATP with $m$ identical agents, targets $\Phi^* = \{s_1,\ldots,s_n\}$, deterministic kernels with $\mathcal{K}^\alpha_{\mathrm{proof}}((1,\varnothing)\mid\mathcal{L},s_j,\tau) = \mathbf{1}[\tau\ge p_j]$, and $\delta^*(s_j) = \varnothing$ for all $j$. Every resolved sequent is automatically fully resolved, and the MATP makespan equals the scheduling makespan.

(2) Under PPI, $o_t$ determines $s_t$, so the POMDP reduces to an MDP. The decision problem for finite-horizon MDPs with finite state and action spaces lies in PSPACE: one can evaluate a given policy in polynomial space by simulating the Markov chain step-by-step, and optimize over policies by enumeration (each policy is a sequence of $H$ action maps, testable in polynomial space).

(3) The component $\mathcal{RS}^0_t$ ranges over all subsets of $\mathcal{T}^*$ (since any subset can arise as the resolved set), giving $|\mathsf{S}| \ge 2^{|\mathcal{T}^*|}$.
\end{proof}

\begin{remark}[Reduction from stochastic online scheduling]
The Non-Conjecturing PPI MATP can also be viewed as an instance of stochastic online scheduling with unrelated machines, preemption, and latent job spawning. Concretely, agents correspond to machines, sequents to jobs, proof kernels to processing time distributions, and dependency revelation to online job arrival. This connection provides access to the scheduling-theory literature for deriving additional approximation and inapproximability results, though the MATP's redundancy (multiple agents on the same sequent) and non-transferable progress (agent-specific cumulative effort) go beyond the standard scheduling model.
\end{remark}

\begin{remark}[Conjecturing and online scheduling]
Enabling conjecturing ($\mathcal{C}_0 \subsetneq \mathcal{F}$) introduces an endogenous online component: the planner generates new sequents during execution via conjecture actions, expanding the task space $\mathcal{CS}_t$. This is qualitatively harder than the non-conjecturing case, as the planner must balance exploration (conjecturing useful lemmas) against exploitation (proving known tasks). The problem becomes a form of stochastic online scheduling with endogenous job generation, for which general-purpose approximation guarantees are scarce. We leave the complexity of the conjecturing case as an open problem and focus on the non-conjecturing case for the remainder of this section.
\end{remark}


\subsection{An Algorithm for the Non-Conjecturing PPI MATP}\label{sec:algorithm}

We now give a concrete policy for the Non-Conjecturing PPI MATP and analyze its performance relative to the lower bounds of Theorem~\ref{thm:lower-bounds}.


\subsubsection{Reactive List Scheduling (RLS)}

The natural algorithmic paradigm is \emph{reactive scheduling}: maintain a queue of known required tasks (initialized to $\Phi^*$, expanded as dependencies are revealed), and apply a greedy heuristic.

\begin{definition}[Reactive List Scheduling]\label{def:rls}
Fix a priority function $\sigma:\mathcal{CS}_0\to\mathbb{R}$ (lower = higher priority). The \emph{Reactive List Scheduling} (RLS) policy:
\begin{enumerate}
    \item Initialize the \emph{known task queue} $\mathcal{K}_0 = \Phi^*$.
    \item At each timestep $t$: assign each idle agent $\alpha$ to the highest-priority unresolved task in $\mathcal{K}_t$ not currently being processed by another agent, with time budget $\tau=1$; idle if no such task exists.
    \item When task $s$ is resolved at time $t$, observe $\delta^*(s)$ and update:
    \[
    \mathcal{K}_{t+1} \leftarrow \mathcal{K}_t \;\cup\; \{s' : s'\text{ induced by }\psi\in\delta^*(s),\; s'\notin\mathcal{RS}^0_t\}.
    \]
\end{enumerate}
\end{definition}

\begin{definition}[Redundant RLS]\label{def:or-rls}
The \emph{Redundant} variant modifies step (2): if $|\mathcal{K}_t \setminus \mathcal{RS}^0_t| < m$ (fewer unresolved known tasks than agents), excess idle agents are assigned to the highest-priority unresolved task already being processed, enabling parallel attempts.
\end{definition}

RLS is a fully specified, polynomial-time policy that requires no knowledge of $\delta^*$ beyond what has been revealed. The only design choice is the priority function $\sigma$; we show below that the approximation guarantee holds for \emph{any} choice of $\sigma$.


\subsubsection{Approximation Guarantee}

\begin{theorem}[RLS approximation]\label{thm:rls-approx}
Consider a Non-Conjecturing PPI MATP instance with $m$ identical agents and single-assignment RLS (Definition~\ref{def:rls}). Then for each realization $\omega$:
\[
C_{\max}^{\mathrm{RLS}}(\omega) \;\le\; \frac{W(\omega)}{m} + \left(1 - \frac{1}{m}\right) D(\omega),
\]
and consequently:
\[
\mathbb{E}[C_{\max}^{\mathrm{RLS}}] \;\le\; \left(2 - \frac{1}{m}\right)\,\mathbb{E}[\mathrm{OPT}].
\]
\end{theorem}

\begin{proof}
Fix a realization $\omega$. Let $s^*$ be the last task in $\mathcal{T}^*$ to complete under RLS. Partition the interval $[0, C_{s^*}]$ into \emph{busy} timesteps (all $m$ agents active) and \emph{starved} timesteps (at least one agent idle).

\emph{Busy time.} At each busy timestep, $m$ units of work are performed. The total work on tasks in $\mathcal{T}^*$ is at most $W(\omega) = \sum_{s\in\mathcal{T}^*} T_s(\omega)$ (identical agents). Thus the total busy time satisfies $B(\omega) \le W(\omega)/m$.

\emph{Starved time.} At a starved timestep, every unresolved task in $\mathcal{K}_t$ is already being processed by some agent, and the remaining agents are idle. Trace the revelation chain that terminates at $s^*$: let $s^*_\ell \to \cdots \to s^*_1 \to s^*_0 = s^*$ be the path in $G^{\mathrm{rev}}$ from a root task $s^*_\ell\in\Phi^*$ to $s^*$. Each starved interval is attributable to waiting for some task in this chain to complete and reveal the next link. The chain tasks' processing times contribute both to the busy time (via $W/m$) and to the starved time. Separating contributions:

The total time decomposes as:
\[
C_{\max}^{\mathrm{RLS}}(\omega) \le \underbrace{\frac{W(\omega) - D_{\mathrm{chain}}(\omega)}{m}}_{\text{busy time net of chain}} + \underbrace{D_{\mathrm{chain}}(\omega)}_{\text{chain time}} = \frac{W(\omega)}{m} + \left(1-\frac{1}{m}\right)D_{\mathrm{chain}}(\omega),
\]
where $D_{\mathrm{chain}}(\omega) = \sum_{i=0}^\ell T_{s^*_i}(\omega) \le D(\omega)$.

Since $\mathrm{OPT}(\omega) \ge \max\{W(\omega)/m,\, D(\omega)\}$ by Theorem~\ref{thm:lower-bounds}:
\[
C_{\max}^{\mathrm{RLS}}(\omega) \le \frac{W(\omega)}{m} + \left(1-\frac{1}{m}\right)D(\omega) \le \left(2-\frac{1}{m}\right)\mathrm{OPT}(\omega).
\]
Taking expectations completes the proof.
\end{proof}

\begin{remark}[Analogy with Graham's list scheduling]
Theorem~\ref{thm:rls-approx} is the MATP analogue of Graham's classical $(2-1/m)$-approximation for $P\mid\mathrm{prec}\mid C_{\max}$. The proof structure is identical: decompose time into busy and idle periods, bound busy time by $W/m$, bound idle time by the longest chain. The key difference is that Graham's ``longest precedence chain'' (bounding forced idleness due to \emph{operational} precedence) is replaced by the serial depth (bounding forced idleness due to \emph{informational} precedence---the latent revelation of dependencies).
\end{remark}


\subsubsection{Redundancy on Revelation Chains}

When the revelation DAG is deep (long chains) but narrow redundancy provides significant speedup.

\begin{proposition}[Redundancy speedup]\label{prop:or-benefit}
Consider a Non-Conjecturing PPI MATP with $m$ identical agents, $\Phi^* = \{s_0\}$, and a pure revelation chain $s_0\to s_1\to\cdots\to s_k$ with i.i.d.\ $\mathrm{Exp}(\lambda)$ processing times. Then:
\begin{enumerate}
    \item Single-assignment RLS: $\mathbb{E}[C_{\max}] = (k+1)/\lambda$.
    \item Redundant RLS: $\mathbb{E}[C_{\max}] = (k+1)/(m\lambda)$.
\end{enumerate}
\end{proposition}

\begin{proof}
The chain is processed serially (each task is discovered upon completion of the previous). Under single-assignment, each task takes $\mathrm{Exp}(\lambda)$ time. Under a redundant policy, all $m$ agents work on each task; by the memoryless property, the minimum of $m$ independent $\mathrm{Exp}(\lambda)$ variables is $\mathrm{Exp}(m\lambda)$.
\end{proof}

\begin{remark}[Task parallelism vs.\ Redundant parallelism]
Redundancy is maximally beneficial on pure revelation chains (where task parallelism is impossible) and provides no benefit when $|\mathcal{K}_t| \ge m$ (enough independent tasks for full utilization). The optimal policy must balance these two modes based on the DAG structure and processing time distributions---a tradeoff that is NP-hard to optimize in general.
\end{remark}


% \subsubsection{Summary}

% \begin{table}[ht]
% \centering
% \renewcommand{\arraystretch}{1.3}
% \begin{tabular}{lll}
% \toprule
% \textbf{Result} & \textbf{Bound} & \textbf{Reference} \\
% \midrule
% Lower bound (any policy) & $\max\{W/m,\;D\}$ & Thm.~\ref{thm:lower-bounds} \\[2pt]
% Upper bound (RLS, identical) & $(2-\frac{1}{m})\cdot\mathrm{OPT}$ & Thm.~\ref{thm:rls-approx} \\[2pt]
% Hardness (decision) & NP-hard (even $\delta^*\equiv\varnothing$) & Thm.~\ref{thm:complexity} \\[2pt]
% Membership (PPI) & PSPACE (MDP) & Thm.~\ref{thm:complexity} \\[2pt]
% Redundancy (chains) & $m\times$ speedup & Prop.~\ref{prop:or-benefit} \\
% \bottomrule
% \end{tabular}
% \caption{Complexity and approximation landscape for the Non-Conjecturing PPI MATP with expected makespan objective.}
% \label{tab:complexity}
% \end{table}


\subsection{Limits of Centralization}\label{sec:limits-centralization}

The centralized MATP provides a strong baseline: a single planner allocates theorem-proving effort across agents to optimize a global objective. Theorem~\ref{thm:ppi-relaxation} shows that PPI is without loss for the planner, and the RLS algorithm achieves a $(2-1/m)$-approximation to the optimal makespan under PPI. However, the centralized formulation presumes a level of control that is unrealistic in the motivating multiagent setting.

In practice, each prover controls its own compute, possesses private search state that is not fully communicable, and decides independently whether to share intermediate results. Once these assumptions are relaxed, the problem is no longer one of single-agent control but of \emph{strategic interaction} under partial information: agents have individual incentives (credit, reputation, competitive advantage) that may conflict with the social objective of minimizing makespan.

This raises two fundamental questions that the centralized formulation cannot address:
\begin{enumerate}
    \item If agents are self-interested, will they voluntarily prove and publish results that benefit others? Or will they withhold private information, duplicate effort, and underinvest in shared infrastructure?
    
    \item Can decentralized agents achieve makespan comparable to the centralized optimum, and if so, through what mechanism?
\end{enumerate}

We address these questions in the next chapter by formulating the problem as a partially observable stochastic game (POSG), identifying the publication externality as the core market failure, and introducing a market mechanism that internalizes this externality and provides bounded efficiency guarantees.